%==============================================================================
%  Highlights — Elsevier **별도 제출 파일**
%
%  원고(PDF 앞쪽 별쪽)와 별개로 투고 시스템에 올리는 항목이다.
%  ★ 손으로 고치지 말 것 — tools/make_highlights.py 가 원고에서 생성한다.
%    원고를 고친 뒤 그 스크립트를 다시 돌리면 여기도 따라온다.
%  컴파일: xelatex highlights.tex
%==============================================================================
\documentclass[11pt]{article}
\usepackage[a4paper, margin=25mm]{geometry}
\usepackage{kotex}
\usepackage{fontspec}
\setmainfont{TeX Gyre Termes}[Ligatures=TeX]
\setmainhangulfont{HCR Batang}[AutoFakeBold=2.0, AutoFakeSlant=0.2]
\usepackage{enumitem}
\pagestyle{empty}
\begin{document}

\noindent{\Large\bfseries Highlights}

\vspace{4pt}
\begin{itemize}[leftmargin=1.2em, itemsep=3pt, topsep=4pt]
\item Two judgment layers share one 25 s period; only the slow one is decoupled
\item Pointing at score-map pixels keeps the action space invariant to swarm size
\item Critic-free group relative optimization with counterfactual credit assignment
\item Across 720 paired episodes the two layers are additive in 9 of 10 metrics
\item Commander skill is explained by plan stability, not assignment coverage
\end{itemize}

\vspace{18pt}
\noindent{\large\bfseries 연구 하이라이트}\;{\small (국문 원고 기준 --- 참고용)}

\vspace{4pt}
\begin{itemize}[leftmargin=1.2em, itemsep=3pt, topsep=4pt]
\item 같은 25초 주기의 두 판단 계층을 지연 차이만으로 비동기 분리한 방어 체계
\item 점수맵 픽셀 지목으로 행동공간을 격자에 가둔, 적선 수에 불변인 협동 정책
\item 가치망 없는 그룹 상대 최적화와 반사실적 크레딧으로 다중 방어정 협동 학습
\item 720 에피소드 대응표본 분해 결과 두 계층의 기여는 10지표 중 9지표에서 가산적
\item 지휘관 성능은 배정 커버리지가 아니라 재계획 간 계획 안정성이 설명한다
\end{itemize}

\end{document}
